\documentclass{article}
\usepackage{amsmath,amssymb,amsthm}
\usepackage{algorithm}
\usepackage{algpseudocode}
\usepackage{bm}
\usepackage{hyperref}
\newtheorem{theorem}{Theorem}
\newtheorem{lemma}{Lemma}
\newtheorem{assumption}{Assumption}
\newcommand{\EE}{\mathbb{E}}
\newcommand{\RR}{\mathbb{R}}
\newcommand{\norm}[1]{\left\lVert#1\right\rVert}
\newcommand{\inner}[2]{\left\langle #1,#2\right\rangle}
\newcommand{\prox}{\operatorname{prox}}
\begin{document}

\section*{Algorithm Name}
\textbf{Differentially Private Stochastic Gradient Descent with Gaussian Noise (DP‑SGD)}  

The algorithm performs standard SGD updates on a convex, $L$‑smooth loss while adding i.i.d. Gaussian noise to each stochastic gradient to guarantee $(\varepsilon,\delta)$‑differential privacy.

\section*{Mathematical Formulation}
Let $f:\RR^{d}\rightarrow\RR$ be the empirical risk
\[
f(w)=\frac{1}{n}\sum_{i=1}^{n}\ell(w;z_i),\qquad 
\ell(\cdot;z_i)\ \text{convex and }L\text{-smooth}.
\]
At iteration $t$ a minibatch $\mathcal{B}_t\subset\{1,\dots ,n\}$ of size $b$ is sampled uniformly without replacement.  
Define the (clipped) stochastic gradient
\[
g_t = \frac{1}{b}\sum_{i\in\mathcal{B}_t}\operatorname{clip}\!\big(\nabla\ell(w_t;z_i),C\big),
\]
where $\operatorname{clip}(v,C)=v\min\{1,\frac{C}{\|v\|}\}$ bounds the $\ell_2$ norm by $C>0$.  
Add Gaussian noise
\[
\tilde g_t = g_t + \xi_t,\qquad \xi_t\sim\mathcal N(0,\sigma^2\mathbf I_d).
\]
The DP‑SGD update is
\[
w_{t+1}= w_t - \eta_t \tilde g_t,
\tag{1}
\]
with step size $\eta_t>0$.  
The noise scale $\sigma$ is chosen according to the moments accountant (or R\'enyi DP) to satisfy a total privacy budget $(\varepsilon,\delta)$ after $T$ iterations.

\section*{Pseudocode}
\begin{algorithm}[h]
\caption{DP‑SGD}
\begin{algorithmic}[1]
\State \textbf{Input:} data $\{z_i\}_{i=1}^n$, loss $\ell$, clipping norm $C$, noise scale $\sigma$, minibatch size $b$, stepsizes $\{\eta_t\}_{t=0}^{T-1}$, iterations $T$.
\State Initialize $w_0\in\RR^{d}$.
\For{$t=0$ \textbf{to} $T-1$}
    \State Sample a minibatch $\mathcal B_t$ of size $b$ uniformly.
    \State Compute per‑example gradients $g_{t,i}= \nabla\ell(w_t;z_i),\; i\in\mathcal B_t$.
    \State Clip: $\tilde g_{t,i}= \operatorname{clip}(g_{t,i},C)$.
    \State Average: $g_t=\frac{1}{b}\sum_{i\in\mathcal B_t}\tilde g_{t,i}$.
    \State Sample noise $\xi_t\sim\mathcal N(0,\sigma^2\mathbf I_d)$.
    \State Noisy gradient: $\tilde g_t = g_t+\xi_t$.
    \State Update: $w_{t+1}= w_t-\eta_t\tilde g_t$.
\EndFor
\State \textbf{Output:} $w_T$ (or average $\bar w_T=\frac{1}{T}\sum_{t=0}^{T-1}w_t$).
\end{algorithmic}
\end{algorithm}

\section*{Dry Run Example}
Consider the one‑dimensional quadratic $f(w)=(w-3)^2$, thus $\nabla f(w)=2(w-3)$.  
Set $w_0=0$, constant stepsize $\eta=0.1$, clipping norm $C=+\infty$ (no clipping), noise variance $\sigma^2=0$ (deterministic for illustration), minibatch size $b=1$, and $T=3$.

\begin{center}
\begin{tabular}{c|c|c|c}
$t$ & $g_t=\nabla f(w_t)$ & $w_{t+1}=w_t-\eta g_t$ & $f(w_{t+1})$\\ \hline
0 & $2(0-3)=-6$ & $0-0.1(-6)=0.6$ & $(0.6-3)^2=5.76$\\
1 & $2(0.6-3)=-4.8$ & $0.6-0.1(-4.8)=1.08$ & $(1.08-3)^2=3.6864$\\
2 & $2(1.08-3)=-3.84$ & $1.08-0.1(-3.84)=1.464$ & $(1.464-3)^2=2.3689$\\
\end{tabular}
\end{center}
The objective value decreases at each iteration, illustrating the basic behaviour of (1) without noise.

\newpage
\section*{Assumptions}
\begin{assumption}[Convexity]
\label{as:convex}
For all $w,w'\in\RR^d$, $f(w')\ge f(w)+\inner{\nabla f(w)}{w'-w}$.
\end{assumption}
\begin{assumption}[Smoothness]
\label{as:smooth}
$f$ is $L$‑smooth: $\forall w,w',\;
\|\nabla f(w)-\nabla f(w')\|\le L\|w-w'\|$.
Equivalently,
\[
f(w')\le f(w)+\inner{\nabla f(w)}{w'-w}+\frac{L}{2}\|w'-w\|^2 .
\]
\end{assumption}
\begin{assumption}[Bounded Gradient after Clipping]
\label{as:clip}
After clipping, $\|g_t\|\le C$ almost surely.
\end{assumption}
\begin{assumption}[Noise Distribution]
\label{as:noise}
$\xi_t\sim\mathcal N(0,\sigma^2\mathbf I_d)$ independent of all past iterates and data.
\end{assumption}
\begin{assumption}[Stepsizes]
\label{as:steps}
The stepsizes are non‑increasing and satisfy $\eta_t\le \frac{1}{L}$ for all $t$.
\end{assumption}

\section*{Main Convergence Theorem}
\begin{theorem}[Expected Suboptimality of DP‑SGD]\label{thm:main}
Under Assumptions \ref{as:convex}–\ref{as:steps}, let $w^*$ be a minimizer of $f$ and define $D:=\|w_0-w^*\|$. For any $T\ge1$ the iterate average $\bar w_T:=\frac{1}{T}\sum_{t=0}^{T-1}w_t$ satisfies
\[
\EE\!\big[f(\bar w_T)-f(w^*)\big]\;\le\;
\frac{D^2}{2\eta T}
\;+\;
\frac{\eta C^2}{2}
\;+\;
\frac{\eta\sigma^2 d}{2}.
\tag{2}
\]
Choosing a constant stepsize $\eta=\frac{D}{C\sqrt{T}}$ (ignoring the $\sigma$ term) yields
\[
\EE\!\big[f(\bar w_T)-f(w^*)\big]=O\!\left(\frac{C D}{\sqrt{T}}+\frac{\sigma\sqrt{d}\,D}{\sqrt{T}}\right).
\]
\end{theorem}

\section*{Theoretical Properties}
\begin{itemize}
\item \textbf{Stability.} The clipping operation guarantees that the sensitivity of each gradient is bounded by $C/b$, which together with Gaussian noise yields a finite privacy loss.
\item \textbf{Hyperparameter Sensitivity.} Larger clipping $C$ reduces bias from clipping but increases required noise $\sigma$ for a fixed privacy budget, degrading the $\sigma^2 d$ term in (2). The stepsize $\eta$ trades off the deterministic term $D^2/(2\eta T)$ against the variance terms $\eta C^2/2$ and $\eta\sigma^2 d/2$.
\item \textbf{Comparison to non‑private SGD.} Setting $\sigma=0$ and removing clipping recovers the classic $O(D^2/(\eta T)+\eta G^2/2)$ bound for SGD with bounded gradients $G$, matching standard results.
\item \textbf{Effect of Dimensionality.} The noise contribution scales linearly with $d$, reflecting the curse of dimensionality typical for DP mechanisms.
\end{itemize}

\section*{Discussion}
The bound (2) shows that DP‑SGD attains the optimal $O(1/\sqrt{T})$ convergence rate for convex smooth objectives up to an additive term proportional to the noise variance. By calibrating $\sigma$ via the moments accountant, one can allocate a privacy budget while still guaranteeing sublinear decay of the expected suboptimality. The theorem also highlights the necessity of clipping: without a bound on $\|g_t\|$, the Gaussian mechanism would either violate privacy or require infinite noise.

\newpage
\section*{Preliminaries (Definitions and Lemmas)}
\begin{lemma}[Three‑point identity]\label{lem:3pt}
For any $a,b\in\RR^d$ and any $u\in\RR^d$,
\[
\|a-u\|^2 = \|b-u\|^2 + 2\inner{a-b}{b-u} + \|a-b\|^2 .
\]
\end{lemma}
\begin{proof}
Expand $\|a-u\|^2 = \|b+(a-b)-u\|^2$ and collect terms. ∎
\end{proof}

\begin{lemma}[Bound on noisy gradient norm]\label{lem:noisy}
Under Assumptions \ref{as:clip} and \ref{as:noise},
\[
\EE\big[\|\tilde g_t\|^2\mid \mathcal F_t\big]\le C^2+\sigma^2 d,
\]
where $\mathcal F_t$ denotes the sigma‑algebra generated by $\{w_0,\dots,w_t\}$.
\end{lemma}
\begin{proof}
Conditioning on $\mathcal F_t$, $g_t$ is deterministic with $\|g_t\|\le C$, and $\xi_t$ is independent with $\EE\|\xi_t\|^2=\sigma^2 d$. Then
\[
\EE\|\tilde g_t\|^2 = \EE\|g_t+\xi_t\|^2 = \|g_t\|^2 + \EE\|\xi_t\|^2 \le C^2+\sigma^2 d .
\qquad\qquad\qquad\qquad\qquad\qquad\qquad\qquad\qquad\qquad\;\; \blacksquare
\]
\end{proof}

\section*{Main Proof (Step‑by‑Step Derivation)}
We prove Theorem \ref{thm:main}. Let $\mathcal F_t$ be as above.

\noindent\textbf{[Step 1]} Expand the squared distance after one update using Lemma \ref{lem:3pt} with $a=w_{t+1}$, $b=w_t$, $u=w^*$:
\[
\|w_{t+1}-w^*\|^2 = \|w_t-w^*\|^2 -2\eta_t\inner{\tilde g_t}{w_t-w^*} + \eta_t^2\|\tilde g_t\|^2 .
\tag{3}
\]

\noindent\textbf{[Step 2]} Rearrange (3) to isolate the inner product term:
\[
2\eta_t\inner{\tilde g_t}{w_t-w^*}
= \|w_t-w^*\|^2-\|w_{t+1}-w^*\|^2 + \eta_t^2\|\tilde g_t\|^2 .
\tag{4}
\]

\noindent\textbf{[Step 3]} Take conditional expectation $\EE[\cdot\mid\mathcal F_t]$ on both sides. Using linearity and the fact that $w_t$ and $w^*$ are $\mathcal F_t$‑measurable,
\[
2\eta_t\EE\!\big[\inner{\tilde g_t}{w_t-w^*}\mid\mathcal F_t\big]
= \|w_t-w^*\|^2-\EE\!\big[\|w_{t+1}-w^*\|^2\mid\mathcal F_t\big] + \eta_t^2\EE\!\big[\|\tilde g_t\|^2\mid\mathcal F_t\big].
\tag{5}
\]

\noindent\textbf{[Step 4]} Decompose $\tilde g_t=g_t+\xi_t$ and use unbiasedness of the noise ($\EE[\xi_t]=0$) to obtain
\[
\EE\!\big[\inner{\tilde g_t}{w_t-w^*}\mid\mathcal F_t\big]
= \inner{g_t}{w_t-w^*}.
\tag{6}
\]

\noindent\textbf{[Step 5]} Apply convexity (Assumption \ref{as:convex}) to $g_t$ (which is the gradient of the clipped empirical loss on the minibatch). Since clipping does not change the direction of the subgradient, we have
\[
\inner{g_t}{w_t-w^*} \ge f(w_t)-f(w^*).
\tag{7}
\]

\noindent\textbf{[Step 6]} Substitute (6) and (7) into (5):
\[
2\eta_t\big(f(w_t)-f(w^*)\big)
\le \|w_t-w^*\|^2-\EE\!\big[\|w_{t+1}-w^*\|^2\mid\mathcal F_t\big] + \eta_t^2\EE\!\big[\|\tilde g_t\|^2\mid\mathcal F_t\big].
\tag{8}
\]

\noindent\textbf{[Step 7]} Use Lemma \ref{lem:noisy} to bound the last term:
\[
\EE\!\big[\|\tilde g_t\|^2\mid\mathcal F_t\big]\le C^2+\sigma^2 d .
\tag{9}
\]

\noindent\textbf{[Step 8]} Insert (9) into (8) and divide by $2\eta_t$:
\[
f(w_t)-f(w^*) \le
\frac{\|w_t-w^*\|^2-\EE\!\big[\|w_{t+1}-w^*\|^2\mid\mathcal F_t\big]}{2\eta_t}
+\frac{\eta_t}{2}\big(C^2+\sigma^2 d\big).
\tag{10}
\]

\noindent\textbf{[Step 9]} Take total expectation $\EE[\cdot]$ on both sides of (10). The telescoping sum over $t=0,\dots,T-1$ yields
\[
\sum_{t=0}^{T-1}\EE\!\big[f(w_t)-f(w^*)\big]
\le
\frac{\|w_0-w^*\|^2}{2\eta_{\min}}
+\frac{1}{2}\sum_{t=0}^{T-1}\eta_t\big(C^2+\sigma^2 d\big),
\tag{11}
\]
where $\eta_{\min}:=\min_{t}\eta_t$ (by Assumption \ref{as:steps} all $\eta_t\le 1/L$, thus the denominator is well defined).

\noindent\textbf{[Step 10]} Divide (11) by $T$ and use convexity of $f$ to replace the average of function values by the function evaluated at the average iterate:
\[
\EE\!\big[f(\bar w_T)-f(w^*)\big]\le
\frac{\|w_0-w^*\|^2}{2\eta_{\min}T}
+\frac{1}{2T}\sum_{t=0}^{T-1}\eta_t\big(C^2+\sigma^2 d\big).
\tag{12}
\]

\noindent\textbf{[Step 11]} For a constant stepsize $\eta_t=\eta$ the sum simplifies:
\[
\EE\!\big[f(\bar w_T)-f(w^*)\big]\le
\frac{D^2}{2\eta T}
+\frac{\eta}{2}\big(C^2+\sigma^2 d\big),
\]
with $D:=\|w_0-w^*\|$, which is exactly bound (2). ∎

\section*{Rate Derivation (With Constants)}
From (12) with constant $\eta$:
\[
\EE\!\big[f(\bar w_T)-f(w^*)\big]\le
\underbrace{\frac{D^2}{2\eta T}}_{\text{deterministic bias}}
\;+\;
\underbrace{\frac{\eta C^2}{2}}_{\text{clipping variance}}
\;+\;
\underbrace{\frac{\eta\sigma^2 d}{2}}_{\text{privacy noise}} .
\tag{13}
\]
Optimizing $\eta$ for the first two terms (ignoring the $\sigma$ term) gives $\eta^\star = \frac{D}{C\sqrt{T}}$, leading to
\[
\EE\!\big[f(\bar w_T)-f(w^*)\big]\le
\frac{DC}{\sqrt{T}}+\frac{D\sigma\sqrt{d}}{\sqrt{T}} .
\tag{14}
\]
Thus DP‑SGD attains the optimal $O(1/\sqrt{T})$ rate with an additional factor $\sigma\sqrt{d}$ stemming from the privacy noise.

\section*{Special Cases}
\begin{itemize}
\item \textbf{No privacy noise ($\sigma=0$).} The bound reduces to the classic SGD rate $\frac{D^2}{2\eta T}+\frac{\eta C^2}{2}$.
\item \textbf{Exact gradients (minibatch $b=n$ and no clipping).} Here $C$ can be taken as $G:=\max_w\|\nabla f(w)\|$, and the bound matches standard results for full‑batch gradient descent with Gaussian perturbation.
\item \textbf{High‑dimensional limit.} If $d\to\infty$ while keeping $\sigma^2 d$ constant (e.g., $\sigma\propto 1/\sqrt{d}$), the noise term remains bounded and the rate does not degrade with dimension.
\end{itemize}

\end{document}